% Línea temporal de la integración europea, 1945-2026 (recupera, rehecha, la
% «imagen 2» del tema antiguo; Víctor, 16/11/2026). Eje vertical de años; a la
% izquierda, los tratados (firma y entrada en vigor) y los hitos
% institucionales; a la derecha, las ampliaciones (con las banderas de los
% Estados que entran, dibujadas en banderas.tikz) y los hitos económicos. Rojo:
% fracasos, crisis y la retirada británica. Datos a 16/11/2026.
% Autor: Víctor Gutiérrez Marcos
\documentclass[tikz,border=4pt]{standalone}
% ==========================================================================
% tikz-tcee.tex — Estilos comunes de los gráficos TikZ del temario
% Autor: Víctor Gutiérrez Marcos
%
% Cada gráfico es un fichero standalone en fuentes/<tema>/graficos/:
%
%   \documentclass[tikz,border=4pt]{standalone}
%   % ==========================================================================
% tikz-tcee.tex — Estilos comunes de los gráficos TikZ del temario
% Autor: Víctor Gutiérrez Marcos
%
% Cada gráfico es un fichero standalone en fuentes/<tema>/graficos/:
%
%   \documentclass[tikz,border=4pt]{standalone}
%   % ==========================================================================
% tikz-tcee.tex — Estilos comunes de los gráficos TikZ del temario
% Autor: Víctor Gutiérrez Marcos
%
% Cada gráfico es un fichero standalone en fuentes/<tema>/graficos/:
%
%   \documentclass[tikz,border=4pt]{standalone}
%   \input{../../../latex/tikz-tcee}
%   \begin{document}
%   \begin{tikzpicture}[tcee]
%     \ejes{Q}{P}{6}{5}
%     \draw[curva1] (0.5,4.5) -- (5,0.8) node[etiqueta,right] {$D$};
%     \capa{2}{\draw[curva2] ...;}   % en el vídeo aparece en el paso 2
%   \end{tikzpicture}
%   \end{document}
%
% Paleta: la de la web (morado, dorado, salvia) más tres colores de apoyo
% distinguibles en blanco y negro por el trazo.
% ==========================================================================
\usepackage[T1]{fontenc}
\usepackage{newpxtext,newpxmath}
\usepackage{xcolor}
\usepackage{tikz,pgfplots}
\pgfplotsset{compat=1.18}
\usetikzlibrary{arrows.meta,calc,positioning,patterns,decorations.pathreplacing,intersections,shapes.misc,backgrounds}

\definecolor{tceemorado}{HTML}{5F2987}
\definecolor{tceedorado}{HTML}{B8860B}
\definecolor{tceeverde}{HTML}{3C7D3A}
\definecolor{tceeazul}{HTML}{1F5F99}
\definecolor{tceerojo}{HTML}{B22222}
\definecolor{tceenaranja}{HTML}{E07B28}
\definecolor{tceegris}{HTML}{767171}
\definecolor{tceesalvia}{HTML}{E2EFD9}

% Capas para el vídeo: \capa{n}{…} dibuja su contenido solo a partir del paso n.
% En el PDF, la web y Word se ve todo; generar-video.py compila el gráfico una
% vez por paso con \def\tceepaso{k} para que cada elemento aparezca al narrarlo.
% Vectores con flecha encima, como en el texto (tcee.sty): \vect{x}, \vect{p}
\providecommand{\vect}[1]{\vec{#1}}
\providecommand{\tceepaso}{1000}
\newcommand{\capa}[2]{\ifnum#1>\tceepaso\relax\else#2\fi}
% Gráficos por pasos del mismo tamaño (propuesta del 3.A.29, 11/11/2026): con
% \def\tceefantasma{1} antes del \input (p. ej., x-paso1.tex:
% «\def\tceepaso{3}\def\tceefantasma{1}\input{x.tex}»), las capas posteriores
% al paso se dibujan invisibles y ocupan su sitio. Sin \tceefantasma (el vídeo),
% no cambia nada.
\ifdefined\tceefantasma
  \renewcommand{\capa}[2]{\ifnum#1>\tceepaso\relax\begin{scope}[opacity=0]#2\end{scope}\else#2\fi}
\fi

\tikzset{
  tcee/.style={line cap=round,line join=round,font=\small,>={Stealth[length=2.2mm]}},
  eje/.style={->,thick,black},
  curva1/.style={very thick,tceemorado},
  curva2/.style={very thick,tceeazul},
  curva3/.style={very thick,tceeverde},
  curva4/.style={very thick,tceedorado},
  curvanueva/.style={very thick,tceerojo},
  desplazada/.style={very thick,dashed},
  guia/.style={thin,dashed,tceegris},
  punto/.style={circle,fill=black,inner sep=1.3pt},
  etiqueta/.style={font=\small,inner sep=2pt},
  flecha/.style={->,thick,tceerojo},
  area/.style={fill=tceemorado,fill opacity=0.18,draw=none},
  area2/.style={fill=tceedorado,fill opacity=0.22,draw=none},
}

% \ejes{etiqueta x}{etiqueta y}{largo x}{alto y}
\newcommand{\ejes}[4]{%
  \draw[eje] (0,0) -- (#3,0) node[below left=2pt and 0pt] {#1};
  \draw[eje] (0,0) -- (0,#4) node[left] {#2};
  \node[below left] at (0,0) {$0$};}
% \proyeccion{(x,y)}{etq x}{etq y}: líneas guía a los ejes con etiquetas
\newcommand{\proyeccion}[3]{%
  \draw[guia] let \p1=#1 in (\x1,0) node[below] {#2} |- (0,\y1) node[left] {#3};}

% --------------------------------------------------------------------------
% Paleta «dictamen» (dic…): gráficos de datos de la Guía del dictamen (1D01) y
% de cualquier gráfico de coyuntura que quiera los mismos colores. Un concepto,
% un color en todos los gráficos (CRITERIOS §7, «tres reglas», 24/10/2026);
% tabla completa en fuentes/1D01/_trabajo/paleta-1D01.md. No cambia nada de
% lo anterior: son colores y estilos nuevos.
% Actividad: el PIB (y sus variantes: PIB por persona, PIB mundial), morado.
\definecolor{dicpib}{HTML}{5F2987}
% Demanda interna, gama cálida: consumo (naranjas: privado oscuro, público
% claro), inversión (verde) y existencias (gris con trama, aparte).
\definecolor{dicconspriv}{HTML}{C76A1F}
\definecolor{dicconspub}{HTML}{F5C690}
\definecolor{dicfbcf}{HTML}{2A6B3C}
\definecolor{dicexistencias}{HTML}{B4B4B4}
% Demanda externa y comercio, gama fría (azul claro, que se separa del morado)
\definecolor{dicexterna}{HTML}{6AA6DA}
% Zonas: zona del euro (azul marino), España (rojo), resto de países (grises)
\definecolor{dicze}{HTML}{17406D}
\definecolor{dicespana}{HTML}{C8382E}
\definecolor{dicotro}{HTML}{8A8A8A}
\definecolor{dicotroclaro}{HTML}{B5B5B5}
\definecolor{dicotromasclaro}{HTML}{DADADA}
% Precios: inflación general (negro, continuo) y subyacente (negro,
% discontinuo); componentes: energía (bermellón), alimentos (amarillo) y los
% de la subyacente en azules (bienes industriales claro, servicios oscuro).
\definecolor{dicinflacion}{HTML}{1A1A1A}
\definecolor{dicenergia}{HTML}{D55E00}
\definecolor{dicalimentos}{HTML}{F2D45C}
\definecolor{dicbine}{HTML}{7FB2E0}
\definecolor{dicservicios}{HTML}{1F4E89}
% Trabajo: empleo (verde azulado), horas (el mismo, claro), tasa de actividad
% (verde azulado oscuro), paro (teja); productividad (dorado de la web).
\definecolor{dicempleo}{HTML}{1B7F6B}
\definecolor{dichoras}{HTML}{86CBB8}
\definecolor{dicactividad}{HTML}{0F5247}
\definecolor{dicparo}{HTML}{A23B2A}
\definecolor{dicproductividad}{HTML}{B8860B}
% Hogares: renta disponible (pizarra) y tasa de ahorro (negro); el consumo es
% dicconspriv.
\definecolor{dicrenta}{HTML}{4A5A6A}
\definecolor{dicahorro}{HTML}{1A1A1A}
% Política monetaria: tipos del BCE en azules (depósito, el de la zona del
% euro; operaciones principales, medio; facilidad marginal, claro); balance:
% préstamos a los bancos (dorado claro), resto del activo (gris claro).
\definecolor{dicdeposito}{HTML}{17406D}
\definecolor{dicopf}{HTML}{4F86C6}
\definecolor{dicfmc}{HTML}{9EC5EA}
\definecolor{dicprestamos}{HTML}{D9A520}
\definecolor{dicresto}{HTML}{CFCFCF}
% Ciclo: series adelantadas, coincidentes (el PIB) y retrasadas
\definecolor{dicadelantada}{HTML}{D55E00}
\definecolor{dicretrasada}{HTML}{6AA6DA}
% Referencias: rejilla, previsión, sombreados
\definecolor{dicreferencia}{HTML}{767171}

\tikzset{
  % barra de un componente: \addplot[dicbarra=dicfbcf]
  dicbarra/.style={fill=#1, draw=#1!65!black, line width=0.3pt},
  % existencias: gris con trama, para que se distinga también en gris
  dicbarratrama/.style={fill=#1, draw=#1!65!black, line width=0.3pt,
    postaction={pattern=north east lines, pattern color=white}},
  % línea de una serie; con halo blanco cuando pasa sobre barras
  diclinea/.style={very thick, #1},
  diclineahalo/.style={very thick, #1, preaction={draw=white, line width=3.4pt, solid}},
  % puntos de una serie (capa aparte, después de la línea)
  dicpuntos/.style={only marks, mark=*, mark size=1.8pt,
    mark options={fill=white, draw=#1, line width=0.9pt, solid}},
  % línea del cero o de un umbral, después de las barras
  diccero/.style={thin, black},
  dicumbral/.style={thin, dashed, dicreferencia},
}
% Capas de los gráficos de datos (CRITERIOS §7, «tres reglas»), de abajo
% arriba: cuadrícula → sombreados (dicfondo) → barras y áreas (dicbarras) → ejes, cero y umbrales
% (axis lines) → líneas (main) → puntos (dicpuntos) → rótulos (axis
% descriptions). No basta el orden del código: pgfplots dibuja cuadrícula,
% ejes y marcas en sus propias capas (y area style pone los ejes encima), y
% dentro de un axis los \draw, \fill y \node no admiten capa (van con los
% datos de main); por eso todo lo que se dibuja dentro del axis es un \addplot.
% Uso: \begin{axis}[…, dicapas] (al final de las opciones);
%   barras y áreas: \addplot[on layer=dicbarras, …];
%   sombreados:     \dicsombra[estilo]{x1}{y1}{x2}{y2} (rectángulo, en dicfondo,
%                   bajo las barras y sobre la cuadrícula);
%   cero y umbrales: \dicraya[estilo]{(x1,y1) (x2,y2)} (en axis lines);
%   líneas: \addplot normal (main); puntos: \addplot[on layer=dicpuntos, …]
%   (las marcas de cualquier serie van solas a dicpuntos);
%   rótulos: \dicrotulo[opciones del nodo]{(x,y)}{texto} (en axis descriptions),
%   siempre sin fondo (nada de fill=white: si chocan con una línea, se desplazan;
%   CRITERIOS §7, regla 4).
% Los \capa{n} del vídeo no cambian: siguen envolviendo instrucciones completas.
\pgfdeclarelayer{dicfondo}
\pgfdeclarelayer{dicbarras}
\pgfdeclarelayer{dicpuntos}
\pgfplotsset{
  % (el segundo argumento reparte cuadrícula, ejes y rótulos como el estándar)
  layers/dic/.define layer set={axis background, axis grid, dicfondo, dicbarras,
    axis ticks, axis lines, axis tick labels, pre main, main, dicpuntos,
    axis descriptions, axis foreground}{/pgfplots/layers/standard},
  % leyenda sin caja blanca: toma el fondo del gráfico, también si el gráfico va
  % dentro de un cuadro gris (CRITERIOS §7, regla 4, 26/10/2026)
  dicapas/.style={axis on top=false, set layers=dic, mark layer=dicpuntos,
    legend style={fill=none}},
  dicsuelto/.style={forget plot, stack plots=false, sharp plot, mark=none},
}
\newcommand{\dicraya}[2][]{%
  \addplot[on layer=axis lines, dicsuelto, diccero, #1] coordinates {#2};}
\newcommand{\dicsombra}[5][]{%
  \addplot[on layer=dicfondo, dicsuelto, draw=none, #1]
    coordinates {(#2,#3) (#4,#3) (#4,#5) (#2,#5)};}
\newcommand{\dicrotulo}[3][]{%
  \addplot[on layer=axis descriptions, dicsuelto, draw=none] coordinates {#2}
    node[pos=0, #1] {#3};}

%   \begin{document}
%   \begin{tikzpicture}[tcee]
%     \ejes{Q}{P}{6}{5}
%     \draw[curva1] (0.5,4.5) -- (5,0.8) node[etiqueta,right] {$D$};
%     \capa{2}{\draw[curva2] ...;}   % en el vídeo aparece en el paso 2
%   \end{tikzpicture}
%   \end{document}
%
% Paleta: la de la web (morado, dorado, salvia) más tres colores de apoyo
% distinguibles en blanco y negro por el trazo.
% ==========================================================================
\usepackage[T1]{fontenc}
\usepackage{newpxtext,newpxmath}
\usepackage{xcolor}
\usepackage{tikz,pgfplots}
\pgfplotsset{compat=1.18}
\usetikzlibrary{arrows.meta,calc,positioning,patterns,decorations.pathreplacing,intersections,shapes.misc,backgrounds}

\definecolor{tceemorado}{HTML}{5F2987}
\definecolor{tceedorado}{HTML}{B8860B}
\definecolor{tceeverde}{HTML}{3C7D3A}
\definecolor{tceeazul}{HTML}{1F5F99}
\definecolor{tceerojo}{HTML}{B22222}
\definecolor{tceenaranja}{HTML}{E07B28}
\definecolor{tceegris}{HTML}{767171}
\definecolor{tceesalvia}{HTML}{E2EFD9}

% Capas para el vídeo: \capa{n}{…} dibuja su contenido solo a partir del paso n.
% En el PDF, la web y Word se ve todo; generar-video.py compila el gráfico una
% vez por paso con \def\tceepaso{k} para que cada elemento aparezca al narrarlo.
% Vectores con flecha encima, como en el texto (tcee.sty): \vect{x}, \vect{p}
\providecommand{\vect}[1]{\vec{#1}}
\providecommand{\tceepaso}{1000}
\newcommand{\capa}[2]{\ifnum#1>\tceepaso\relax\else#2\fi}
% Gráficos por pasos del mismo tamaño (propuesta del 3.A.29, 11/11/2026): con
% \def\tceefantasma{1} antes del \input (p. ej., x-paso1.tex:
% «\def\tceepaso{3}\def\tceefantasma{1}\input{x.tex}»), las capas posteriores
% al paso se dibujan invisibles y ocupan su sitio. Sin \tceefantasma (el vídeo),
% no cambia nada.
\ifdefined\tceefantasma
  \renewcommand{\capa}[2]{\ifnum#1>\tceepaso\relax\begin{scope}[opacity=0]#2\end{scope}\else#2\fi}
\fi

\tikzset{
  tcee/.style={line cap=round,line join=round,font=\small,>={Stealth[length=2.2mm]}},
  eje/.style={->,thick,black},
  curva1/.style={very thick,tceemorado},
  curva2/.style={very thick,tceeazul},
  curva3/.style={very thick,tceeverde},
  curva4/.style={very thick,tceedorado},
  curvanueva/.style={very thick,tceerojo},
  desplazada/.style={very thick,dashed},
  guia/.style={thin,dashed,tceegris},
  punto/.style={circle,fill=black,inner sep=1.3pt},
  etiqueta/.style={font=\small,inner sep=2pt},
  flecha/.style={->,thick,tceerojo},
  area/.style={fill=tceemorado,fill opacity=0.18,draw=none},
  area2/.style={fill=tceedorado,fill opacity=0.22,draw=none},
}

% \ejes{etiqueta x}{etiqueta y}{largo x}{alto y}
\newcommand{\ejes}[4]{%
  \draw[eje] (0,0) -- (#3,0) node[below left=2pt and 0pt] {#1};
  \draw[eje] (0,0) -- (0,#4) node[left] {#2};
  \node[below left] at (0,0) {$0$};}
% \proyeccion{(x,y)}{etq x}{etq y}: líneas guía a los ejes con etiquetas
\newcommand{\proyeccion}[3]{%
  \draw[guia] let \p1=#1 in (\x1,0) node[below] {#2} |- (0,\y1) node[left] {#3};}

% --------------------------------------------------------------------------
% Paleta «dictamen» (dic…): gráficos de datos de la Guía del dictamen (1D01) y
% de cualquier gráfico de coyuntura que quiera los mismos colores. Un concepto,
% un color en todos los gráficos (CRITERIOS §7, «tres reglas», 24/10/2026);
% tabla completa en fuentes/1D01/_trabajo/paleta-1D01.md. No cambia nada de
% lo anterior: son colores y estilos nuevos.
% Actividad: el PIB (y sus variantes: PIB por persona, PIB mundial), morado.
\definecolor{dicpib}{HTML}{5F2987}
% Demanda interna, gama cálida: consumo (naranjas: privado oscuro, público
% claro), inversión (verde) y existencias (gris con trama, aparte).
\definecolor{dicconspriv}{HTML}{C76A1F}
\definecolor{dicconspub}{HTML}{F5C690}
\definecolor{dicfbcf}{HTML}{2A6B3C}
\definecolor{dicexistencias}{HTML}{B4B4B4}
% Demanda externa y comercio, gama fría (azul claro, que se separa del morado)
\definecolor{dicexterna}{HTML}{6AA6DA}
% Zonas: zona del euro (azul marino), España (rojo), resto de países (grises)
\definecolor{dicze}{HTML}{17406D}
\definecolor{dicespana}{HTML}{C8382E}
\definecolor{dicotro}{HTML}{8A8A8A}
\definecolor{dicotroclaro}{HTML}{B5B5B5}
\definecolor{dicotromasclaro}{HTML}{DADADA}
% Precios: inflación general (negro, continuo) y subyacente (negro,
% discontinuo); componentes: energía (bermellón), alimentos (amarillo) y los
% de la subyacente en azules (bienes industriales claro, servicios oscuro).
\definecolor{dicinflacion}{HTML}{1A1A1A}
\definecolor{dicenergia}{HTML}{D55E00}
\definecolor{dicalimentos}{HTML}{F2D45C}
\definecolor{dicbine}{HTML}{7FB2E0}
\definecolor{dicservicios}{HTML}{1F4E89}
% Trabajo: empleo (verde azulado), horas (el mismo, claro), tasa de actividad
% (verde azulado oscuro), paro (teja); productividad (dorado de la web).
\definecolor{dicempleo}{HTML}{1B7F6B}
\definecolor{dichoras}{HTML}{86CBB8}
\definecolor{dicactividad}{HTML}{0F5247}
\definecolor{dicparo}{HTML}{A23B2A}
\definecolor{dicproductividad}{HTML}{B8860B}
% Hogares: renta disponible (pizarra) y tasa de ahorro (negro); el consumo es
% dicconspriv.
\definecolor{dicrenta}{HTML}{4A5A6A}
\definecolor{dicahorro}{HTML}{1A1A1A}
% Política monetaria: tipos del BCE en azules (depósito, el de la zona del
% euro; operaciones principales, medio; facilidad marginal, claro); balance:
% préstamos a los bancos (dorado claro), resto del activo (gris claro).
\definecolor{dicdeposito}{HTML}{17406D}
\definecolor{dicopf}{HTML}{4F86C6}
\definecolor{dicfmc}{HTML}{9EC5EA}
\definecolor{dicprestamos}{HTML}{D9A520}
\definecolor{dicresto}{HTML}{CFCFCF}
% Ciclo: series adelantadas, coincidentes (el PIB) y retrasadas
\definecolor{dicadelantada}{HTML}{D55E00}
\definecolor{dicretrasada}{HTML}{6AA6DA}
% Referencias: rejilla, previsión, sombreados
\definecolor{dicreferencia}{HTML}{767171}

\tikzset{
  % barra de un componente: \addplot[dicbarra=dicfbcf]
  dicbarra/.style={fill=#1, draw=#1!65!black, line width=0.3pt},
  % existencias: gris con trama, para que se distinga también en gris
  dicbarratrama/.style={fill=#1, draw=#1!65!black, line width=0.3pt,
    postaction={pattern=north east lines, pattern color=white}},
  % línea de una serie; con halo blanco cuando pasa sobre barras
  diclinea/.style={very thick, #1},
  diclineahalo/.style={very thick, #1, preaction={draw=white, line width=3.4pt, solid}},
  % puntos de una serie (capa aparte, después de la línea)
  dicpuntos/.style={only marks, mark=*, mark size=1.8pt,
    mark options={fill=white, draw=#1, line width=0.9pt, solid}},
  % línea del cero o de un umbral, después de las barras
  diccero/.style={thin, black},
  dicumbral/.style={thin, dashed, dicreferencia},
}
% Capas de los gráficos de datos (CRITERIOS §7, «tres reglas»), de abajo
% arriba: cuadrícula → sombreados (dicfondo) → barras y áreas (dicbarras) → ejes, cero y umbrales
% (axis lines) → líneas (main) → puntos (dicpuntos) → rótulos (axis
% descriptions). No basta el orden del código: pgfplots dibuja cuadrícula,
% ejes y marcas en sus propias capas (y area style pone los ejes encima), y
% dentro de un axis los \draw, \fill y \node no admiten capa (van con los
% datos de main); por eso todo lo que se dibuja dentro del axis es un \addplot.
% Uso: \begin{axis}[…, dicapas] (al final de las opciones);
%   barras y áreas: \addplot[on layer=dicbarras, …];
%   sombreados:     \dicsombra[estilo]{x1}{y1}{x2}{y2} (rectángulo, en dicfondo,
%                   bajo las barras y sobre la cuadrícula);
%   cero y umbrales: \dicraya[estilo]{(x1,y1) (x2,y2)} (en axis lines);
%   líneas: \addplot normal (main); puntos: \addplot[on layer=dicpuntos, …]
%   (las marcas de cualquier serie van solas a dicpuntos);
%   rótulos: \dicrotulo[opciones del nodo]{(x,y)}{texto} (en axis descriptions),
%   siempre sin fondo (nada de fill=white: si chocan con una línea, se desplazan;
%   CRITERIOS §7, regla 4).
% Los \capa{n} del vídeo no cambian: siguen envolviendo instrucciones completas.
\pgfdeclarelayer{dicfondo}
\pgfdeclarelayer{dicbarras}
\pgfdeclarelayer{dicpuntos}
\pgfplotsset{
  % (el segundo argumento reparte cuadrícula, ejes y rótulos como el estándar)
  layers/dic/.define layer set={axis background, axis grid, dicfondo, dicbarras,
    axis ticks, axis lines, axis tick labels, pre main, main, dicpuntos,
    axis descriptions, axis foreground}{/pgfplots/layers/standard},
  % leyenda sin caja blanca: toma el fondo del gráfico, también si el gráfico va
  % dentro de un cuadro gris (CRITERIOS §7, regla 4, 26/10/2026)
  dicapas/.style={axis on top=false, set layers=dic, mark layer=dicpuntos,
    legend style={fill=none}},
  dicsuelto/.style={forget plot, stack plots=false, sharp plot, mark=none},
}
\newcommand{\dicraya}[2][]{%
  \addplot[on layer=axis lines, dicsuelto, diccero, #1] coordinates {#2};}
\newcommand{\dicsombra}[5][]{%
  \addplot[on layer=dicfondo, dicsuelto, draw=none, #1]
    coordinates {(#2,#3) (#4,#3) (#4,#5) (#2,#5)};}
\newcommand{\dicrotulo}[3][]{%
  \addplot[on layer=axis descriptions, dicsuelto, draw=none] coordinates {#2}
    node[pos=0, #1] {#3};}

%   \begin{document}
%   \begin{tikzpicture}[tcee]
%     \ejes{Q}{P}{6}{5}
%     \draw[curva1] (0.5,4.5) -- (5,0.8) node[etiqueta,right] {$D$};
%     \capa{2}{\draw[curva2] ...;}   % en el vídeo aparece en el paso 2
%   \end{tikzpicture}
%   \end{document}
%
% Paleta: la de la web (morado, dorado, salvia) más tres colores de apoyo
% distinguibles en blanco y negro por el trazo.
% ==========================================================================
\usepackage[T1]{fontenc}
\usepackage{newpxtext,newpxmath}
\usepackage{xcolor}
\usepackage{tikz,pgfplots}
\pgfplotsset{compat=1.18}
\usetikzlibrary{arrows.meta,calc,positioning,patterns,decorations.pathreplacing,intersections,shapes.misc,backgrounds}

\definecolor{tceemorado}{HTML}{5F2987}
\definecolor{tceedorado}{HTML}{B8860B}
\definecolor{tceeverde}{HTML}{3C7D3A}
\definecolor{tceeazul}{HTML}{1F5F99}
\definecolor{tceerojo}{HTML}{B22222}
\definecolor{tceenaranja}{HTML}{E07B28}
\definecolor{tceegris}{HTML}{767171}
\definecolor{tceesalvia}{HTML}{E2EFD9}

% Capas para el vídeo: \capa{n}{…} dibuja su contenido solo a partir del paso n.
% En el PDF, la web y Word se ve todo; generar-video.py compila el gráfico una
% vez por paso con \def\tceepaso{k} para que cada elemento aparezca al narrarlo.
% Vectores con flecha encima, como en el texto (tcee.sty): \vect{x}, \vect{p}
\providecommand{\vect}[1]{\vec{#1}}
\providecommand{\tceepaso}{1000}
\newcommand{\capa}[2]{\ifnum#1>\tceepaso\relax\else#2\fi}
% Gráficos por pasos del mismo tamaño (propuesta del 3.A.29, 11/11/2026): con
% \def\tceefantasma{1} antes del \input (p. ej., x-paso1.tex:
% «\def\tceepaso{3}\def\tceefantasma{1}\input{x.tex}»), las capas posteriores
% al paso se dibujan invisibles y ocupan su sitio. Sin \tceefantasma (el vídeo),
% no cambia nada.
\ifdefined\tceefantasma
  \renewcommand{\capa}[2]{\ifnum#1>\tceepaso\relax\begin{scope}[opacity=0]#2\end{scope}\else#2\fi}
\fi

\tikzset{
  tcee/.style={line cap=round,line join=round,font=\small,>={Stealth[length=2.2mm]}},
  eje/.style={->,thick,black},
  curva1/.style={very thick,tceemorado},
  curva2/.style={very thick,tceeazul},
  curva3/.style={very thick,tceeverde},
  curva4/.style={very thick,tceedorado},
  curvanueva/.style={very thick,tceerojo},
  desplazada/.style={very thick,dashed},
  guia/.style={thin,dashed,tceegris},
  punto/.style={circle,fill=black,inner sep=1.3pt},
  etiqueta/.style={font=\small,inner sep=2pt},
  flecha/.style={->,thick,tceerojo},
  area/.style={fill=tceemorado,fill opacity=0.18,draw=none},
  area2/.style={fill=tceedorado,fill opacity=0.22,draw=none},
}

% \ejes{etiqueta x}{etiqueta y}{largo x}{alto y}
\newcommand{\ejes}[4]{%
  \draw[eje] (0,0) -- (#3,0) node[below left=2pt and 0pt] {#1};
  \draw[eje] (0,0) -- (0,#4) node[left] {#2};
  \node[below left] at (0,0) {$0$};}
% \proyeccion{(x,y)}{etq x}{etq y}: líneas guía a los ejes con etiquetas
\newcommand{\proyeccion}[3]{%
  \draw[guia] let \p1=#1 in (\x1,0) node[below] {#2} |- (0,\y1) node[left] {#3};}

% --------------------------------------------------------------------------
% Paleta «dictamen» (dic…): gráficos de datos de la Guía del dictamen (1D01) y
% de cualquier gráfico de coyuntura que quiera los mismos colores. Un concepto,
% un color en todos los gráficos (CRITERIOS §7, «tres reglas», 24/10/2026);
% tabla completa en fuentes/1D01/_trabajo/paleta-1D01.md. No cambia nada de
% lo anterior: son colores y estilos nuevos.
% Actividad: el PIB (y sus variantes: PIB por persona, PIB mundial), morado.
\definecolor{dicpib}{HTML}{5F2987}
% Demanda interna, gama cálida: consumo (naranjas: privado oscuro, público
% claro), inversión (verde) y existencias (gris con trama, aparte).
\definecolor{dicconspriv}{HTML}{C76A1F}
\definecolor{dicconspub}{HTML}{F5C690}
\definecolor{dicfbcf}{HTML}{2A6B3C}
\definecolor{dicexistencias}{HTML}{B4B4B4}
% Demanda externa y comercio, gama fría (azul claro, que se separa del morado)
\definecolor{dicexterna}{HTML}{6AA6DA}
% Zonas: zona del euro (azul marino), España (rojo), resto de países (grises)
\definecolor{dicze}{HTML}{17406D}
\definecolor{dicespana}{HTML}{C8382E}
\definecolor{dicotro}{HTML}{8A8A8A}
\definecolor{dicotroclaro}{HTML}{B5B5B5}
\definecolor{dicotromasclaro}{HTML}{DADADA}
% Precios: inflación general (negro, continuo) y subyacente (negro,
% discontinuo); componentes: energía (bermellón), alimentos (amarillo) y los
% de la subyacente en azules (bienes industriales claro, servicios oscuro).
\definecolor{dicinflacion}{HTML}{1A1A1A}
\definecolor{dicenergia}{HTML}{D55E00}
\definecolor{dicalimentos}{HTML}{F2D45C}
\definecolor{dicbine}{HTML}{7FB2E0}
\definecolor{dicservicios}{HTML}{1F4E89}
% Trabajo: empleo (verde azulado), horas (el mismo, claro), tasa de actividad
% (verde azulado oscuro), paro (teja); productividad (dorado de la web).
\definecolor{dicempleo}{HTML}{1B7F6B}
\definecolor{dichoras}{HTML}{86CBB8}
\definecolor{dicactividad}{HTML}{0F5247}
\definecolor{dicparo}{HTML}{A23B2A}
\definecolor{dicproductividad}{HTML}{B8860B}
% Hogares: renta disponible (pizarra) y tasa de ahorro (negro); el consumo es
% dicconspriv.
\definecolor{dicrenta}{HTML}{4A5A6A}
\definecolor{dicahorro}{HTML}{1A1A1A}
% Política monetaria: tipos del BCE en azules (depósito, el de la zona del
% euro; operaciones principales, medio; facilidad marginal, claro); balance:
% préstamos a los bancos (dorado claro), resto del activo (gris claro).
\definecolor{dicdeposito}{HTML}{17406D}
\definecolor{dicopf}{HTML}{4F86C6}
\definecolor{dicfmc}{HTML}{9EC5EA}
\definecolor{dicprestamos}{HTML}{D9A520}
\definecolor{dicresto}{HTML}{CFCFCF}
% Ciclo: series adelantadas, coincidentes (el PIB) y retrasadas
\definecolor{dicadelantada}{HTML}{D55E00}
\definecolor{dicretrasada}{HTML}{6AA6DA}
% Referencias: rejilla, previsión, sombreados
\definecolor{dicreferencia}{HTML}{767171}

\tikzset{
  % barra de un componente: \addplot[dicbarra=dicfbcf]
  dicbarra/.style={fill=#1, draw=#1!65!black, line width=0.3pt},
  % existencias: gris con trama, para que se distinga también en gris
  dicbarratrama/.style={fill=#1, draw=#1!65!black, line width=0.3pt,
    postaction={pattern=north east lines, pattern color=white}},
  % línea de una serie; con halo blanco cuando pasa sobre barras
  diclinea/.style={very thick, #1},
  diclineahalo/.style={very thick, #1, preaction={draw=white, line width=3.4pt, solid}},
  % puntos de una serie (capa aparte, después de la línea)
  dicpuntos/.style={only marks, mark=*, mark size=1.8pt,
    mark options={fill=white, draw=#1, line width=0.9pt, solid}},
  % línea del cero o de un umbral, después de las barras
  diccero/.style={thin, black},
  dicumbral/.style={thin, dashed, dicreferencia},
}
% Capas de los gráficos de datos (CRITERIOS §7, «tres reglas»), de abajo
% arriba: cuadrícula → sombreados (dicfondo) → barras y áreas (dicbarras) → ejes, cero y umbrales
% (axis lines) → líneas (main) → puntos (dicpuntos) → rótulos (axis
% descriptions). No basta el orden del código: pgfplots dibuja cuadrícula,
% ejes y marcas en sus propias capas (y area style pone los ejes encima), y
% dentro de un axis los \draw, \fill y \node no admiten capa (van con los
% datos de main); por eso todo lo que se dibuja dentro del axis es un \addplot.
% Uso: \begin{axis}[…, dicapas] (al final de las opciones);
%   barras y áreas: \addplot[on layer=dicbarras, …];
%   sombreados:     \dicsombra[estilo]{x1}{y1}{x2}{y2} (rectángulo, en dicfondo,
%                   bajo las barras y sobre la cuadrícula);
%   cero y umbrales: \dicraya[estilo]{(x1,y1) (x2,y2)} (en axis lines);
%   líneas: \addplot normal (main); puntos: \addplot[on layer=dicpuntos, …]
%   (las marcas de cualquier serie van solas a dicpuntos);
%   rótulos: \dicrotulo[opciones del nodo]{(x,y)}{texto} (en axis descriptions),
%   siempre sin fondo (nada de fill=white: si chocan con una línea, se desplazan;
%   CRITERIOS §7, regla 4).
% Los \capa{n} del vídeo no cambian: siguen envolviendo instrucciones completas.
\pgfdeclarelayer{dicfondo}
\pgfdeclarelayer{dicbarras}
\pgfdeclarelayer{dicpuntos}
\pgfplotsset{
  % (el segundo argumento reparte cuadrícula, ejes y rótulos como el estándar)
  layers/dic/.define layer set={axis background, axis grid, dicfondo, dicbarras,
    axis ticks, axis lines, axis tick labels, pre main, main, dicpuntos,
    axis descriptions, axis foreground}{/pgfplots/layers/standard},
  % leyenda sin caja blanca: toma el fondo del gráfico, también si el gráfico va
  % dentro de un cuadro gris (CRITERIOS §7, regla 4, 26/10/2026)
  dicapas/.style={axis on top=false, set layers=dic, mark layer=dicpuntos,
    legend style={fill=none}},
  dicsuelto/.style={forget plot, stack plots=false, sharp plot, mark=none},
}
\newcommand{\dicraya}[2][]{%
  \addplot[on layer=axis lines, dicsuelto, diccero, #1] coordinates {#2};}
\newcommand{\dicsombra}[5][]{%
  \addplot[on layer=dicfondo, dicsuelto, draw=none, #1]
    coordinates {(#2,#3) (#4,#3) (#4,#5) (#2,#5)};}
\newcommand{\dicrotulo}[3][]{%
  \addplot[on layer=axis descriptions, dicsuelto, draw=none] coordinates {#2}
    node[pos=0, #1] {#3};}

% Banderas simplificadas de los Estados europeos, dibujadas en TikZ para los
% gráficos del tema 3.B.36 (circulos-europa, linea-temporal). Las banderas
% nacionales son de dominio público; aquí se dibujan a mano, simplificadas
% (escudos y emblemas esquematizados), en proporción 3:2 salvo Suiza y el
% Vaticano (cuadradas). Colores aproximados a los oficiales.
% Uso: \bandera[escala]{(x,y)}{código}: centrada en (x,y); con escala 0.2,
% 0,6 x 0,4 cm. Códigos ISO de dos letras (GR, Grecia; XK, Kosovo; GB, Reino
% Unido; EU, la bandera europea).
% Autor: Víctor Gutiérrez Marcos
\makeatletter
% estrella de 5 puntas: \bandera@estrella{x}{y}{radio}{color}
\newcommand{\bandera@estrella}[4]{%
  \fill[#4] ($(#1,#2)+(90:#3)$)
    \foreach \k in {1,...,9}{ -- ($(#1,#2)+({90+36*\k}:{ifthenelse(mod(\k,2)==1,0.4*#3,#3)})$) } -- cycle;}
% franjas horizontales (de arriba abajo) y verticales (de izquierda a derecha)
\newcommand{\bandera@h}[3]{\fill[#1] (0,4/3) rectangle (3,2);\fill[#2] (0,2/3) rectangle (3,4/3);\fill[#3] (0,0) rectangle (3,2/3);}
\newcommand{\bandera@hdos}[2]{\fill[#1] (0,1) rectangle (3,2);\fill[#2] (0,0) rectangle (3,1);}
\newcommand{\bandera@v}[3]{\fill[#1] (0,0) rectangle (1,2);\fill[#2] (1,0) rectangle (2,2);\fill[#3] (2,0) rectangle (3,2);}
% cruz nórdica: fondo, cruz y, opcional, cruz interior
\newcommand{\bandera@nordica}[3]{\fill[#1] (0,0) rectangle (3,2);
  \fill[#2] (0.8,0) rectangle (1.4,2);\fill[#2] (0,0.7) rectangle (3,1.3);
  \ifx&#3&\else\fill[#3] (0.95,0) rectangle (1.25,2);\fill[#3] (0,0.85) rectangle (3,1.15);\fi}
% escudo pequeño genérico: \bandera@escudo{x}{y}{ancho}{relleno}{borde}
\newcommand{\bandera@escudo}[5]{%
  \filldraw[fill=#4,draw=#5,line width=0.15pt] ($(#1,#2)+(-#3/2,#3*0.55)$) -- ++(#3,0) -- ++(0,-#3*0.6)
    .. controls +(0,-#3*0.35) and +(#3*0.3,0) .. ($(#1,#2)+(0,-#3*0.62)$)
    .. controls +(-#3*0.3,0) and +(0,-#3*0.35) .. ($(#1,#2)+(-#3/2,-#3*0.05)$) -- cycle;}
\definecolor{bnrojo}{HTML}{D52B1E}\definecolor{bnazul}{HTML}{003F87}\definecolor{bnoro}{HTML}{FFCE00}
\definecolor{bnverde}{HTML}{009246}\definecolor{bnazulc}{HTML}{0072CE}
\newcommand{\bandera@AT}{\fill[white] (0,0) rectangle (3,2);\fill[red!85!black] (0,4/3) rectangle (3,2);\fill[red!85!black] (0,0) rectangle (3,2/3);}
\newcommand{\bandera@BE}{\bandera@v{black}{bnoro}{red!90!black}}
\newcommand{\bandera@BG}{\bandera@h{white}{bnverde!90!black}{bnrojo}}
\newcommand{\bandera@HR}{\bandera@h{red}{white}{blue!70!black}%
  \fill[white] (1.15,1.45) rectangle (1.85,0.55);
  \foreach \i in {0,...,4}{\foreach \j in {0,...,4}{\pgfmathtruncatemacro{\bpar}{mod(\i+\j,2)}
    \ifnum\bpar=0 \fill[red] ({1.15+0.14*\i},{1.45-0.14*\j}) rectangle ++(0.14,-0.14);\fi}}
  \draw[blue!70!black,line width=0.2pt] (1.15,1.45) rectangle (1.85,0.75);}
\newcommand{\bandera@EE}{\bandera@h{bnazulc}{black}{white}}
\newcommand{\bandera@FI}{\bandera@nordica{white}{blue!75!black}{}}
\newcommand{\bandera@FR}{\bandera@v{blue!80!black}{white}{bnrojo}}
\newcommand{\bandera@DE}{\bandera@h{black}{red!90!black}{bnoro}}
\newcommand{\bandera@GR}{\fill[white] (0,0) rectangle (3,2);
  \foreach \i in {0,2,4,6,8}{\fill[blue!70!black] (0,{2-\i*2/9}) rectangle (3,{2-(\i+1)*2/9});}
  \fill[blue!70!black] (0,2) rectangle (10/9,{2-10/9});
  \fill[white] (4/9,2) rectangle (6/9,{2-10/9});\fill[white] (0,{2-4/9}) rectangle (10/9,{2-6/9});}
\newcommand{\bandera@IT}{\bandera@v{bnverde}{white}{red!85!black}}
\newcommand{\bandera@LV}{\fill[white] (0,0) rectangle (3,2);\fill[red!55!black] (0,1.2) rectangle (3,2);\fill[red!55!black] (0,0) rectangle (3,0.8);}
\newcommand{\bandera@LT}{\bandera@h{bnoro!90!orange}{green!40!black}{red!80!black}}
\newcommand{\bandera@LU}{\bandera@h{red!90!black}{white}{cyan!70!blue}}
\newcommand{\bandera@MT}{\fill[white] (0,0) rectangle (1.5,2);\fill[red!85!black] (1.5,0) rectangle (3,2);
  \fill[gray!55] (0.3,1.55) rectangle (0.6,1.65);\fill[gray!55] (0.4,1.45) rectangle (0.5,1.75);
  \draw[red!85!black,line width=0.15pt] (0.27,1.42) rectangle (0.63,1.78);}
\newcommand{\bandera@NL}{\bandera@h{red!80!black}{white}{blue!70!black}}
\newcommand{\bandera@PT}{\fill[green!45!black] (0,0) rectangle (1.2,2);\fill[red] (1.2,0) rectangle (3,2);
  \draw[bnoro,line width=0.6pt] (1.2,1) circle (0.42);
  \bandera@escudo{1.2}{1}{0.42}{white}{red}\fill[blue!70!black] (1.15,1.05) rectangle (1.25,0.85);}
\newcommand{\bandera@SK}{\bandera@h{white}{blue!75!black}{red}%
  \bandera@escudo{1.05}{1.05}{0.85}{red}{white}
  \fill[white] (1.0,0.85) rectangle (1.1,1.45);\fill[white] (0.85,1.25) rectangle (1.25,1.33);\fill[white] (0.9,1.07) rectangle (1.2,1.15);
  \fill[blue!75!black] (0.75,0.72) .. controls (0.9,0.95) and (1.2,0.95) .. (1.35,0.72) -- cycle;}
\newcommand{\bandera@SI}{\bandera@h{white}{blue!70!black}{red}%
  \bandera@escudo{0.95}{1.35}{0.55}{blue!70!black}{red}
  \fill[white] (0.73,1.12) -- (0.95,1.45) -- (1.17,1.12) -- cycle;
  \fill[bnoro] (0.95,1.62) circle (0.04);\fill[bnoro] (0.85,1.55) circle (0.04);\fill[bnoro] (1.05,1.55) circle (0.04);}
\newcommand{\bandera@ES}{\fill[red!75!black] (0,0) rectangle (3,2);\fill[bnoro!90!orange] (0,0.5) rectangle (3,1.5);
  \bandera@escudo{0.95}{1.0}{0.42}{red!75!black}{red!40!black}
  \fill[bnoro!80!orange] (0.95,1.0) rectangle (1.16,0.78);\fill[white] (0.74,1.0) rectangle (0.95,1.23);
  \fill[gray!70] (0.62,0.62) rectangle (0.68,1.15);\fill[gray!70] (1.22,0.62) rectangle (1.28,1.15);
  \fill[red!60!black] (0.8,1.25) rectangle (1.1,1.33);}
\newcommand{\bandera@CZ}{\bandera@hdos{white}{red!85!black}\fill[blue!70!black] (0,0) -- (1.5,1) -- (0,2) -- cycle;}
\newcommand{\bandera@DK}{\bandera@nordica{red!85!black}{white}{}}
\newcommand{\bandera@HU}{\bandera@h{red!80!black}{white}{green!45!black}}
\newcommand{\bandera@PL}{\bandera@hdos{white}{red!85!black}}
\newcommand{\bandera@RO}{\bandera@v{blue!75!black}{bnoro!95!orange}{red!85!black}}
\newcommand{\bandera@SE}{\bandera@nordica{cyan!45!blue!80!black}{bnoro}{}}
\newcommand{\bandera@IE}{\bandera@v{green!55!black}{white}{orange!90!red}}
\newcommand{\bandera@CY}{\fill[white] (0,0) rectangle (3,2);
  \fill[orange!75!brown] (0.75,1.15) -- (1.05,1.35) -- (1.55,1.38) -- (2.0,1.55) -- (2.35,1.62) -- (2.05,1.35)
    -- (1.95,1.15) -- (1.7,1.05) -- (1.45,0.92) -- (1.15,0.95) -- (0.9,1.0) -- cycle;
  \draw[green!35!black,line width=0.5pt] (1.15,0.75) .. controls (1.3,0.55) and (1.5,0.55) .. (1.5,0.55);
  \draw[green!35!black,line width=0.5pt] (1.85,0.75) .. controls (1.7,0.55) and (1.5,0.55) .. (1.5,0.55);}
\newcommand{\bandera@IS}{\bandera@nordica{blue!70!black}{white}{red!85!black}}
\newcommand{\bandera@NO}{\bandera@nordica{red!80!black}{white}{blue!60!black}}
\newcommand{\bandera@LI}{\bandera@hdos{blue!75!black}{red!85!black}%
  \fill[bnoro] (0.45,1.35) -- (0.45,1.6) -- (0.6,1.48) -- (0.75,1.7) -- (0.9,1.48) -- (1.05,1.6) -- (1.05,1.35) -- cycle;}
\newcommand{\bandera@CH}{\fill[red] (0.5,0) rectangle (2.5,2);
  \fill[white] (1.35,0.4) rectangle (1.65,1.6);\fill[white] (0.9,0.85) rectangle (2.1,1.15);}
\newcommand{\bandera@TR}{\fill[red!90!black] (0,0) rectangle (3,2);
  \fill[white] (1.05,1) circle (0.5);\fill[red!90!black] (1.18,1) circle (0.4);
  \bandera@estrella{1.85}{1}{0.25}{white}}
\newcommand{\bandera@AD}{\fill[blue!80!black] (0,0) rectangle (0.95,2);\fill[bnoro] (0.95,0) rectangle (2.05,2);\fill[red!85!black] (2.05,0) rectangle (3,2);
  \bandera@escudo{1.5}{1.0}{0.5}{bnoro!80!orange}{red!50!black}
  \fill[red!85!black] (1.53,1.0) rectangle (1.72,0.75);\fill[red!85!black] (1.28,1.27) rectangle (1.47,1.0);}
\newcommand{\bandera@SM}{\bandera@hdos{white}{cyan!55!blue!70}%
  \bandera@escudo{1.5}{1.0}{0.5}{cyan!40}{bnoro!70!brown}
  \fill[white] (1.35,0.9) rectangle (1.42,1.15);\fill[white] (1.47,0.9) rectangle (1.54,1.2);\fill[white] (1.59,0.9) rectangle (1.66,1.15);}
\newcommand{\bandera@MC}{\bandera@hdos{red!85!black}{white}}
\newcommand{\bandera@VA}{\fill[bnoro] (0.5,0) rectangle (1.5,2);\fill[white] (1.5,0) rectangle (2.5,2);
  \draw[bnoro!80!orange,line width=0.7pt] (1.7,0.5) -- (2.3,1.2);\draw[gray!60,line width=0.7pt] (2.3,0.5) -- (1.7,1.2);
  \fill[white] (1.85,1.25) rectangle (2.15,1.6);\draw[bnoro!80!orange,line width=0.3pt] (1.85,1.25) rectangle (2.15,1.6);}
\newcommand{\bandera@ME}{\fill[bnoro!80!orange] (0,0) rectangle (3,2);\fill[red!75!black] (0.12,0.12) rectangle (2.88,1.88);
  \fill[bnoro!80!orange] (1.5,0.55) -- (1.2,0.85) -- (0.95,1.4) -- (1.3,1.15) -- (1.5,1.45) -- (1.7,1.15) -- (2.05,1.4) -- (1.8,0.85) -- cycle;}
\newcommand{\bandera@XK}{\fill[blue!65!black] (0,0) rectangle (3,2);
  \fill[bnoro!70!brown] (1.05,0.5) -- (0.95,0.9) -- (1.2,1.25) -- (1.6,1.28) -- (1.95,1.05) -- (2.0,0.7) -- (1.7,0.45) -- (1.35,0.35) -- cycle;
  \foreach \a in {-25,-15,-5,5,15,25}{\bandera@estrella{{1.5+2.6*sin(\a)}}{{1.6-0.6*(1-cos(\a))*5}}{0.09}{white}}}
\newcommand{\bandera@GB}{\fill[blue!60!black] (0,0) rectangle (3,2);
  \draw[white,line width=1.1pt] (0,0) -- (3,2) (0,2) -- (3,0);\draw[red!85!black,line width=0.4pt] (0,0) -- (3,2) (0,2) -- (3,0);
  \fill[white] (1.25,0) rectangle (1.75,2);\fill[white] (0,0.75) rectangle (3,1.25);
  \fill[red!85!black] (1.35,0) rectangle (1.65,2);\fill[red!85!black] (0,0.85) rectangle (3,1.15);}
\newcommand{\bandera@EU}{\fill[blue!70!black] (0,0) rectangle (3,2);
  \foreach \a in {0,30,...,330}{\bandera@estrella{{1.5+0.62*cos(\a)}}{{1+0.62*sin(\a)}}{0.1}{bnoro}}}
\newcommand{\bandera@UA}{\bandera@hdos{blue!70!cyan!80!black}{bnoro}}
\newcommand{\bandera@MD}{\bandera@v{blue!75!black}{bnoro}{red!85!black}\bandera@escudo{1.5}{1.0}{0.45}{brown!70!orange}{brown!50!black}}
\newcommand{\bandera@AL}{\fill[red!85!black] (0,0) rectangle (3,2);
  \fill[black] (1.5,0.45) -- (1.25,0.8) -- (0.8,1.45) -- (1.3,1.2) -- (1.5,1.55) -- (1.7,1.2) -- (2.2,1.45) -- (1.75,0.8) -- cycle;}
\newcommand{\bandera@RS}{\bandera@h{red!80!black}{blue!60!black}{white}\bandera@escudo{1.05}{1.05}{0.5}{red!80!black}{bnoro}}
\newcommand{\bandera@MK}{\fill[red!85!black] (0,0) rectangle (3,2);
  \foreach \a in {0,45,...,315}{\fill[bnoro] (1.5,1) -- ($(1.5,1)+(\a-7:2)$) -- ($(1.5,1)+(\a+7:2)$) -- cycle;}
  \fill[red!85!black] (0,0) rectangle (3,-1);\fill[red!85!black] (0,2) rectangle (3,3);
  \fill[bnoro] (1.5,1) circle (0.3);\fill[red!85!black] (1.5,1) circle (0.24);\fill[bnoro] (1.5,1) circle (0.2);}
\newcommand{\bandera@BA}{\fill[blue!65!black] (0,0) rectangle (3,2);\fill[bnoro] (0.8,2) -- (2.2,2) -- (2.2,0) -- cycle;
  \foreach \k in {0,...,7}{\bandera@estrella{{0.6+0.2*\k}}{{1.95-0.286*\k}}{0.09}{white}}}
\newcommand{\bandera@GE}{\fill[white] (0,0) rectangle (3,2);\fill[red] (1.3,0) rectangle (1.7,2);\fill[red] (0,0.8) rectangle (3,1.2);
  \foreach \x/\y in {0.65/1.5,2.35/1.5,0.65/0.5,2.35/0.5}{\fill[red] ({\x-0.05},{\y-0.15}) rectangle ({\x+0.05},{\y+0.15});\fill[red] ({\x-0.15},{\y-0.05}) rectangle ({\x+0.15},{\y+0.05});}}
% \bandera[escala]{(x,y)}{código}
\newcommand{\bandera}[3][0.2]{%
  \begin{scope}[shift={#2},scale=#1,shift={(-1.5,-1)}]
    \begin{scope}
      \ifnum\pdfstrcmp{#3}{CH}=0 \clip (0.5,0) rectangle (2.5,2);
      \else\ifnum\pdfstrcmp{#3}{VA}=0 \clip (0.5,0) rectangle (2.5,2);
      \else \clip (0,0) rectangle (3,2);\fi\fi
      \csname bandera@#3\endcsname
    \end{scope}
    \ifnum\pdfstrcmp{#3}{CH}=0 \draw[black!35,line width=0.2pt] (0.5,0) rectangle (2.5,2);
    \else\ifnum\pdfstrcmp{#3}{VA}=0 \draw[black!35,line width=0.2pt] (0.5,0) rectangle (2.5,2);
    \else \draw[black!35,line width=0.2pt] (0,0) rectangle (3,2);\fi\fi
  \end{scope}}
\makeatother

\begin{document}
\newcommand{\nocortar}{\hyphenpenalty=10000 \exhyphenpenalty=10000 }
% posición vertical de un año (1945 arriba)
\newcommand{\anio}[1]{\pgfmathsetmacro{\yA}{-(#1-1945)*0.29}}
\newcommand{\pos}[2]{\anio{#1}\pgfmathsetmacro{\yL}{\yA+#2}}
% \izq{año}{desplazamiento del rótulo}{color}{texto}: hito a la izquierda
\newcommand{\izq}[4]{\pos{#1}{#2}%
  \fill[#3] (8.15,\yA) circle (0.07);
  \draw[#3!70,line width=0.4pt] (8.15,\yA) -- (7.75,\yL) -- (7.55,\yL);
  \node[rotizq,text=#3] at (7.5,\yL) {#4};}
% \der{año}{desplazamiento}{color}{texto}: hito a la derecha
\newcommand{\der}[4]{\pos{#1}{#2}%
  \fill[#3] (8.85,\yA) circle (0.07);
  \draw[#3!70,line width=0.4pt] (8.85,\yA) -- (9.25,\yL) -- (9.45,\yL);
  \node[rotder,text=#3] at (9.5,\yL) {#4};}
% \amp{año}{desplazamiento}{banderas}{texto}: ampliación, con las banderas
\newcommand{\amp}[4]{\pos{#1}{#2}%
  \fill[tceeverde] (8.85,\yA) circle (0.08);
  \draw[tceeverde!70,line width=0.4pt] (8.85,\yA) -- (9.25,\yL) -- (9.45,\yL);
  \foreach \cb [count=\k from 0] in {#3}{\pgfmathsetmacro{\xb}{9.78+0.6*\k}\bandera[0.17]{(\xb,\yL)}{\cb}}
  \node[rotder,text=tceeverde!85!black,font=\footnotesize\bfseries] at (9.5,\yL-0.38) {#4};}
\begin{tikzpicture}[tcee,every node/.append style={execute at begin node=\nocortar},
    rotizq/.style={anchor=east,align=right,font=\footnotesize,inner sep=1pt},
    rotder/.style={anchor=west,align=left,font=\footnotesize,inner sep=1pt}]
  % 1: el eje de los años
  \capa{1}{\fill[tceegris!15] (8.15,0.3) rectangle (8.85,-23.6);
    \foreach \a in {1945,...,2026}{\anio{\a}\draw[tceegris!60,line width=0.3pt] (8.15,\yA) -- (8.27,\yA) (8.73,\yA) -- (8.85,\yA);}
    \foreach \a in {1945,1950,...,2025}{\anio{\a}\node[font=\scriptsize\bfseries,text=tceegris,fill=white,inner sep=1pt,rounded corners=1pt] at (8.5,\yA) {\a};}
    \node[font=\scriptsize\bfseries,text=tceemorado,fill=white,inner sep=1pt,rounded corners=1pt] at (8.5,-23.42) {2026};
    \node[font=\small\bfseries,text=tceemorado,anchor=south east] at (7.6,0.35) {Tratados e instituciones};
    \node[font=\small\bfseries,text=tceeverde!85!black,anchor=south west] at (9.4,0.35) {Ampliaciones y economía};}
  % 2: de la posguerra a los Tratados de Roma (1945-1958)
  \capa{2}{\izq{1948.4}{0.15}{tceeazul}{Congreso de La Haya (1948)}
    \izq{1949.35}{-0.05}{tceeazul}{Consejo de Europa (1949)}
    \izq{1950.35}{-0.25}{tceemorado}{\textbf{Declaración Schuman} (9 de mayo de 1950)}
    \izq{1951.3}{-0.55}{tceemorado}{\textbf{Tratado de París}: CECA (en vigor, 1952)}
    \izq{1954.6}{-0.25}{tceerojo}{Fracasa la Comunidad Europea de Defensa (1954)}
    \izq{1957.2}{-0.15}{tceemorado}{\textbf{Tratados de Roma}: CEE y Euratom (en vigor, 1958)}
    \der{1948.3}{0.1}{tceedorado}{Plan Marshall y OECE (1948)}
    \amp{1952.5}{0.05}{FR,DE,IT,BE,NL,LU}{Los Seis fundan la CECA (1952)}}
  % 3: la expansión de los 60 y la primera ampliación (1958-1973)
  \capa{3}{\izq{1965.3}{0.2}{tceemorado}{\textbf{Tratado de fusión} de los ejecutivos (en vigor, 1967)}
    \izq{1965.5}{-0.35}{tceerojo}{Crisis de la silla vacía (1965) y Compromiso de Luxemburgo (1966)}
    \izq{1970.3}{0.05}{tceemorado}{Tratado de Luxemburgo: recursos propios (1970)}
    \der{1960.3}{0}{tceedorado}{Nace la AELC/EFTA (1960)}
    \der{1968.5}{0.15}{tceedorado}{Unión aduanera completa (1 de julio de 1968)}
    \der{1972.3}{0.12}{tceedorado}{Serpiente monetaria (1972)}
    \amp{1973}{-0.2}{GB,IE,DK}{1.ª ampliación (1973): 9 Estados}}
  % 4: crisis y relanzamiento (1973-1990)
  \capa{4}{\izq{1974.9}{-0.05}{tceeazul}{Nace el Consejo Europeo (1974)}
    \izq{1975.55}{-0.4}{tceemorado}{Tratado de Bruselas: Tribunal de Cuentas (1975)}
    \izq{1979.45}{-0.1}{tceeazul}{Primeras elecciones directas al Parlamento (1979)}
    \izq{1985.45}{0.15}{tceedorado}{Libro Blanco del mercado interior (\textsc{Delors}, 1985)}
    \izq{1986.15}{-0.3}{tceemorado}{\textbf{Acta Única Europea} (en vigor, 1987)}
    \der{1979.2}{0.1}{tceedorado}{Sistema Monetario Europeo (1979)}
    \amp{1981}{-0.15}{GR}{2.ª ampliación (1981): 10}
    \amp{1986}{-0.15}{ES,PT}{3.ª ampliación (1986): 12}
    \der{1990.75}{0.05}{tceeazul}{Caída del Muro (1989) y reunificación alemana (1990)}}
  % 5: de la Comunidad a la Unión (1990-2007)
  \capa{5}{\izq{1992.1}{0.1}{tceemorado}{\textbf{Tratado de Maastricht}: la UE y sus 3 pilares (en vigor, 1993)}
    \izq{1997.8}{0.1}{tceemorado}{\textbf{Tratado de Ámsterdam} (en vigor, 1999)}
    \izq{2001.1}{0.05}{tceemorado}{\textbf{Tratado de Niza} (en vigor, 2003)}
    \izq{2004.8}{0}{tceerojo}{Constitución Europea (2004), rechazada en 2005}
    \izq{2007.95}{-0.1}{tceemorado}{\textbf{Tratado de Lisboa}: TUE, TFUE y Carta (en vigor, 2009)}
    \der{1993}{0.15}{tceedorado}{Mercado interior (1993) y EEE (1994)}
    \amp{1995}{-0.05}{AT,FI,SE}{4.ª ampliación (1995): 15}
    \der{1999}{-0.1}{tceedorado}{Nace el euro (1999; billetes en 2002)}
    \amp{2004.33}{0.35}{CZ,EE,CY,LV,LT}{}\bandera[0.17]{(9.78,\yL-0.42)}{HU}\bandera[0.17]{(10.38,\yL-0.42)}{MT}\bandera[0.17]{(10.98,\yL-0.42)}{PL}\bandera[0.17]{(11.58,\yL-0.42)}{SI}\bandera[0.17]{(12.18,\yL-0.42)}{SK}\node[rotder,text=tceeverde!85!black,font=\footnotesize\bfseries] at (9.5,\yL-0.8) {5.ª ampliación (2004): 25};
    \amp{2007}{-0.35}{BG,RO}{Bulgaria y Rumanía (2007): 27}}
  % 6: crisis, Brexit y ampliación pendiente (2008-2026)
  \capa{6}{\izq{2012.2}{0.1}{tceeazul}{Tratado del MEDE y pacto fiscal (TECG, 2012)}
    \izq{2016.5}{0.05}{tceerojo}{Referéndum del Brexit (2016)}
    \izq{2020.1}{0.05}{tceerojo}{El Reino Unido sale de la Unión (31 de enero de 2020)}
    \izq{2023.9}{0.1}{tceeazul}{El Parlamento propone reformar los Tratados (2023)}
    \izq{2026.75}{-0.05}{tceeazul}{La Comisión propone «Preparar una Unión más amplia» (2026)}
    \der{2010}{-0.55}{tceerojo}{Crisis financiera (2008) y del euro (2010-2012)}
    \amp{2013.5}{0.15}{HR}{Croacia (2013): 28}
    \der{2020.5}{0.45}{tceedorado}{\emph{NextGenerationEU} (2020)}
    \amp{2020.08}{0}{GB}{Brexit (2020): 27}
    \amp{2022.5}{-0.3}{UA,MD}{Ucrania y Moldavia, candidatas (2022)}
    \der{2026.1}{-0.2}{tceedorado}{Bulgaria entra en el euro (2026)}
    \der{2026.5}{-0.65}{tceeverde!85!black}{Capítulos abiertos con Ucrania y Moldavia (2026)}}
\end{tikzpicture}
\end{document}
